\section{A sharper one-dimensional result with a shallow ReLU learner}
\label{sec:univariate-Gaussian}
In one input dimension, a one-hidden-layer ReLU learner gives a much
better width exponent while retaining independent Gaussian weights,
raw likelihood, and raw prediction. This section changes the learner
activation to ReLU and uses exactly one hidden layer:
\[
 z_w(x)=\sum_{j=1}^Nv_j\operatorname{ReLU}(a_jx+b_j),\qquad
 a_j,b_j\sim N(0,\nu),\quad v_j\sim N(0,\nu/N^2),\quad0<\nu\le1.
\]
All coordinates are independent. Assume $X\in[-1,1]$ and
\begin{equation}
 P_X(I)\le\kappa |I|^\alpha
 \quad\hbox{for every interval }I\subseteq[-1,1],
 \qquad\kappa\ge1,\quad0<\alpha\le1.
 \label{eq:one-dimensional-regularity}
\end{equation}
A bounded input density gives $\alpha=1$.

\begin{lemma}[Shallow Gaussian path-variation oracle]
\label{lem:shallow-path-interface}
For a deterministic shallow ReLU implementation $q$, put
\[
 P(q)=\sum_j|v_j|\sqrt{a_j^2+b_j^2}.
\]
There is an observable evidence calibration independent of this
implementation such that, with probability at least $1-\delta$,
simultaneously for all $N$ above its width and all $0<\nu\le1$,
\[
 \Pi_SR(z)\le R(q)+\widetilde O\left(
 (P(q)+1)\sqrt{(R(q)+1)/n}+(P(q)+1)^2/n\right).
 \label{eq:one-dimensional-path-oracle}
\]
Labels may have the same marginal sub-Gaussian envelope as before.
The ordinary raw mean obeys the same bound. Logarithms include
$n,P(q),1/\delta$; constants depend on the fixed label envelope.
\end{lemma}
\begin{proof}
This is the path-variation specialization of the Gaussian free-energy
oracle~\cite{adaptive}; we give the normalization and comparison
calculation to specify its application to the present prior.
Set $\omega=N/\nu$ and
\[
 Q=N^{-1}\sum_j(a_j^2+b_j^2)+N\sum_jv_j^2.
\]
Then $\omega Q\sim\chi^2_{3N}$ and $3N/\omega\le3$.
Cauchy--Schwarz gives $P(w)\le Q/2$. The input vector $(x,1)$ has
norm at most $\sqrt2$, so the output envelope is at most $Q/\sqrt2$,
and the energy-ball Rademacher complexity is at most $cQ/\sqrt n$.
The latter follows from ReLU contraction and absolute convexification
of Euclidean unit linear features. The native pointwise variance is
at most one, so its empirical prior risk under bounded unit labels
is at most two.

Positive rescaling of each comparison neuron preserves its function
and balances its two energy contributions. Zero-padding then gives an
exact width-$N$ mean implementation with $Q=2P(q)$ at every eligible
$N$. Shrink all native Gaussian coordinate variances by $\eta^2$.
For the resulting comparison law,
\[
 (2/\omega)\KL(\rho\|\mu)
       \le2P(q)+6\log(1/\eta),
\]
and independence of the first-layer and readout perturbations gives
\[
 \bigl(\E_\rho\|z_w-q\|_\infty^2\bigr)^{1/2}
       \le4\eta\sqrt{P(q)}+2\eta^2.
\]
Taking $\eta=[64n(P(q)+1)^2]^{-1}$ makes the comparison empirical
squared-loss increment at most $c/n$ for bounded labels.
The free-energy certificate with capacity radius linear in $Q$ has
powers $(v,\rho)=(2,1)$ and comparison resource
$\widetilde O(P(q)+1)$. Its scalar root is
\[
 (\lambda')^2=\theta_n[2F'_S(\lambda')+b_n\lambda'],
 \qquad \theta_n=\widetilde O(1/n),\quad b_n=O(\log(n/\delta)),
\]
with the same bounded-risk fallback for small samples. Apply it to
the normalized observable and labels from
Section~\ref{sec:statistical-interface}, and set
$\lambda=M_n^2\lambda'$. The physical observation variance is still
$n\lambda\nu/N$. The normalized-observable transfer changes only
logarithmic factors. This proves \eqref{eq:one-dimensional-path-oracle}.
\end{proof}

\begin{lemma}[A bounded spline with controlled path variation]
\label{lem:spline-path}
If a continuous $f:[-1,1]\to[-1,1]$ has at most $K$ affine
breakpoints, then for every $0<t\le1$ there is a bounded continuous
piecewise-affine $q$ with at most $2K$ breakpoints such that
\[
 \E|q(X)-f(X)|\le6\kappa Kt^\alpha,\qquad
 \operatorname{Lip}(q)\le2/t,
 \qquad P(q)\le c(K+1)/t.
\]
It has a one-hidden-layer ReLU implementation of width at most $2K+2$.
\end{lemma}
\begin{proof}
Group consecutive breakpoints into maximal clusters whose successive
gaps are at most $3t$. For a cluster running from $a$ to $b$, replace
$f$ on $[\max(-1,a-t),\min(1,b+t)]$ by the chord between its endpoint
values. These replacement intervals are disjoint, with gaps greater
than $t$. A cluster containing $k$ breakpoints has interval length at
most $3kt$. The function remains bounded by one, and differs from $f$
only on these intervals. Equation~\eqref{eq:one-dimensional-regularity}
therefore gives error at most
$2\kappa\sum(3kt)^\alpha\le6\kappa Kt^\alpha$.

Each replacement interval has length at least $t$, so its chord slope
is at most $2/t$. Complementary interior intervals have length greater
than $t$ and are affine. At the boundary, any short complementary
interval belongs to an original affine segment extending at least
$t$ to the first original breakpoint; its slope has the same bound.
Thus all slopes have magnitude at most $2/t$.

If the new breakpoints are $u_1<\cdots<u_r$ and the successive slopes
are $s_0,\ldots,s_r$, use the exact representation
\[
 q(x)=q(-1)\operatorname{ReLU}(1)
     +s_0\operatorname{ReLU}(x+1)
     +\sum_{j=1}^r(s_j-s_{j-1})\operatorname{ReLU}(x-u_j).
\]
Here $r\le2K$, the incoming row norms are at most $\sqrt2$, and each
slope jump is at most $4/t$. This proves the path-variation bound.
The case $K=0$ is a bounded affine function and follows directly.
\end{proof}

\begin{theorem}[A raw $C n^{-1/4}$ oracle under a bounded density]
\label{thm:univariate-Gaussian}
Let $\mathcal G_C$ be all one-hidden-layer univariate ReLU networks of
width at most $C$, with arbitrary finite coefficients and biases.
Under \eqref{eq:one-dimensional-regularity} and the marginal
sub-Gaussian label assumption, the shallow Gaussian learner above
has, with probability at least $1-\delta$, simultaneously for all
$N\ge6C+6$ and all $0<\nu\le1$,
\begin{align*}
 R(\bar z_S)\le\Pi_SR(z)
 \le\inf_{g\in\mathcal G_C}R(g)+\widetilde O\bigl(
 &(C+1)\kappa^{1/(\alpha+1)}n^{-\alpha/[2(\alpha+1)]}\\
 &+(C+1)^2\kappa^{2/(\alpha+1)}n^{-\alpha/(\alpha+1)}\bigr).
\end{align*}
In particular, for a bounded density ($\alpha=1$), the excess is
$\widetilde O(C\sqrt\kappa\,n^{-1/4}+C^2\kappa n^{-1/2})$.
The same single calibration competes with every declared comparison
width; it does not use $C,\kappa,\alpha$ or the original coefficients.
A sufficient fixed-accuracy sample bound is
\[
 n\ge\widetilde O\left(
 \kappa^{2/\alpha}
       ((C+1)/\varepsilon)^{2+2/\alpha}\right).
\]
\end{theorem}
\begin{proof}
Choose a deterministic $n^{-2}$-minimizer $g$ of the comparison risk.
The normalized clipped function
$f=\clip(g/M_n)$ has at most $K=3C+2$ breakpoints: each of the
original $C+1$ affine intervals acquires at most two clipping crossings.
Apply Lemma~\ref{lem:spline-path} at
\[
 t=(\kappa\sqrt n)^{-1/(\alpha+1)}\le1,
\]
and multiply the result by $M_n$. The compiled raw network has width
at most $6C+6$, supremum at most $M_n$, and path variation at most
$cM_n(C+1)/t$. Lemma~\ref{lem:compiler-risk} gives a risk increment
$\widetilde O(\kappa(C+1)t^\alpha)$, plus a negligible label-tail term.
Its risk is also at most $O(M_n^2+\E Y^2)$ independently of $C$.
Substitute this path bound into
\eqref{eq:one-dimensional-path-oracle}. The additional terms are
\[
 \widetilde O\left(\frac{C+1}{t\sqrt n}
                         +\frac{(C+1)^2}{t^2n}\right).
\]
The displayed choice of $t$ balances the two leading contributions
and proves the oracle. Solving its first term for sample size gives
the final bound; its squared term is then smaller. Countably many
declared comparison widths can share one event through a summable
confidence allocation, changing only logarithms.
\end{proof}
